%%%
%
% $Autor: Wings $
% $Date: 2024-10-31 11:15:45Z $
% $File: TensorFlowLiteTool.tex $
% $Version: 1 $
%
%%%


\section{TensorFlow Lite as a Tool}

TensorFlow Lite as a tool primarily refers to the package of utilities and functions provided by TensorFlow Lite to convert, optimize, and prepare TensorFlow models for deployment on edge devices, such as mobile phones, microcontrollers and IoT devices. The primary goal of using TensorFlow Lite as a tool is to make machine learning models smaller, faster, and more efficient without compromising much on accuracy. \cite{TensorFlowLite:2024}

\begin{itemize}
	\item \textbf{Model Conversion}: The core of TensorFlow Lite as a tool is the TensorFlow Lite Converter, refer ~\ref{subsec:model_conversion}. This tool converts a standard TensorFlow model into a \SHELL{.tflite} TensorFlow Lite model format.
	\item \textbf{Model Optimization}: TensorFlow Lite offers several optimization techniques that can be applied during the conversion process or as part of a separate step.
\end{itemize}


\subsection{Relationship to Other TensorFlow Lite functionality}
The use of TensorFlow Lite as a tool is complementary to using TensorFlow Lite as a library path or as a library for Arduino. While the library path provides the necessary runtime environment for deploying models, and the Arduino library allows integration on microcontrollers \cite{TensorFlowLiteMicrocontrollers:2024}, the tools ensure that the models are optimized for these environments. The conversion and optimization steps are typically performed once and the resulting model format \SHELL{.tflite} is then deployed across various platforms and devices. \cite{TensorFlowLiteModelConversion:2024}

In a typical development workflow, a developer would first train a model using TensorFlow, convert and optimize it using TensorFlow Lite tools, and then deploy it using the appropriate TensorFlow Lite library, whether on a mobile device, an Arduino board, or another embedded system. This sequence ensures that the model is not only effective but also efficient and suitable for real-world constraints.


\section{TensorFlow Lite, Version 2.12.0}
\subsection{Introduction}
TensorFlow Lite is a lightweight solution for deploying machine learning models on mobile and embedded devices. It is designed to run TensorFlow models on resource-constrained devices with low latency and high efficiency. \cite{TensorFlowLiteGuide:2024}

\subsection{Installation}

To install TensorFlow Lite, run \SHELL{pip install tensorflow} in command prompt window, refer ~\ref{tf_install}

\subsection{Example - Description}

This example demonstrates how to use TensorFlow Lite to perform image classification on a mobile device. The steps include loading a model, preprocessing input data, running inference, and interpreting the results.

\subsection{Example - Manual}

The manual steps for running a TensorFlow Lite model include:

\begin{enumerate}
	\item Convert the TensorFlow model to TensorFlow Lite format.
	\item Load the TensorFlow Lite model on the device.
	\item Preprocess the input data to match the model’s requirements.
	\item Use the TensorFlow Lite interpreter to run the model.
	\item Postprocess the results to interpret the model’s output.
\end{enumerate}

\subsection{Example - Code}

{
	\captionof{code}{Example Code for running Tensorflow Lite model.}\label{ExampleCode_TensorFlow.ino}
	\ArduinoExternal{}{../../Code/EdgeImpulse/Packages/ExampleCodeTensorFlow.ino}
}




\subsection{Example - Files}

The example files include:

\begin{itemize}
	\item \texttt{model.tflite} - The TensorFlow Lite model file.
	\item \texttt{image.jpg} - The input image file for inference.
	\item \texttt{tflite\_example.py} - The Python script to run the model.
\end{itemize}

\subsection{Futher reading}
For more detailed information, refer to the TensorFlow Lite documentation:
\begin{itemize}
	\item \href{https://www.tensorflow.org/lite/guide}{TensorFlow Lite Guide}.
	\item \href{https://www.tensorflow.org/lite/models}{TensorFlow Lite Models}.
	\item \href{https://www.tensorflow.org/lite/convert}{TensorFlow Lite Conversion}.
	
\end{itemize}


